\documentclass[10pt,a4paper]{amsart}
\usepackage{amsmath,amssymb,amsthm,graphicx}
% For commutative diagrams we use Xy-pic
\usepackage[all]{xy}

% We use 2cell for 2-commutative diagrams.
\xyoption{2cell}
% Theorem environments.
%
\theoremstyle{plain}
\newtheorem{theorem}[subsection]{Theorem}
\newtheorem{proposition}[subsection]{Proposition}
\newtheorem{lemma}[subsection]{Lemma}

\theoremstyle{definition}
\newtheorem{definition}[subsection]{Definition}
\newtheorem{example}[subsection]{Example}
\newtheorem{exercise}[subsection]{Exercise}
\newtheorem{situation}[subsection]{Situation}

\theoremstyle{remark}
\newtheorem{remark}[subsection]{Remark}
\newtheorem{remarks}[subsection]{Remarks}

\numberwithin{equation}{subsection}

% Macros
%
\def\lim{\mathop{\mathrm{lim}}\nolimits}
\def\colim{\mathop{\mathrm{colim}}\nolimits}
\def\Spec{\mathop{\mathrm{Spec}}}
\def\Hom{\mathop{\mathrm{Hom}}\nolimits}
\def\Ext{\mathop{\mathrm{Ext}}\nolimits}
\def\SheafHom{\mathop{\mathcal{H}\!\mathit{om}}\nolimits}
\def\SheafExt{\mathop{\mathcal{E}\!\mathit{xt}}\nolimits}
\def\Sch{\mathit{Sch}}
\def\Mor{\mathop{\mathrm{Mor}}\nolimits}
\def\Ob{\mathop{\mathrm{Ob}}\nolimits}
\def\Sh{\mathop{\mathit{Sh}}\nolimits}
\def\NL{\mathop{N\!L}\nolimits}
\def\CH{\mathop{\mathrm{CH}}\nolimits}
\def\proetale{{pro\text{-}\acute{e}tale}}
\def\etale{{\acute{e}tale}}
\def\QCoh{\mathit{QCoh}}
\def\Ker{\mathop{\mathrm{Ker}}}
\def\Im{\mathop{\mathrm{Im}}}
\def\Coker{\mathop{\mathrm{Coker}}}
\def\Coim{\mathop{\mathrm{Coim}}}

% Boxtimes
%
\DeclareMathSymbol{\boxtimes}{\mathbin}{AMSa}{"02}

%
% Macros for moduli stacks/spaces
%
\def\QCohstack{\mathcal{QC}\!\mathit{oh}}
\def\Cohstack{\mathcal{C}\!\mathit{oh}}
\def\Spacesstack{\mathcal{S}\!\mathit{paces}}
\def\Quotfunctor{\mathrm{Quot}}
\def\Hilbfunctor{\mathrm{Hilb}}
\def\Curvesstack{\mathcal{C}\!\mathit{urves}}
\def\Polarizedstack{\mathcal{P}\!\mathit{olarized}}
\def\Complexesstack{\mathcal{C}\!\mathit{omplexes}}
% \Pic is the operator that assigns to X its picard group, usage \Pic(X)
% \Picardstack_{X/B} denotes the Picard stack of X over B
% \Picardfunctor_{X/B} denotes the Picard functor of X over B
\def\Pic{\mathop{\mathrm{Pic}}\nolimits}
\def\Picardstack{\mathcal{P}\!\mathit{ic}}
\def\Picardfunctor{\mathrm{Pic}}
\def\Deformationcategory{\mathcal{D}\!\mathit{ef}}
\usepackage{cleveref}
\crefname{originallemma}{original lemma}{original lemmas}
\Crefname{originallemma}{Original lemma}{Original lemmas}
\makeatletter
\newlabel{algebra-lemma-filter-Noetherian-module}{{00L0}{source}{}{sourcepointer-1}{}}
\newlabel{algebra-lemma-filter-Noetherian-module@cref}{{[originallemma][0][]00L0}{[1][0][]source}{}{}{}}
\makeatother
\crefname{originallemma}{original lemma}{original lemmas}
\Crefname{originallemma}{Original lemma}{Original lemmas}
\makeatletter
\newlabel{algebra-lemma-filter-minimal-primes-in-support}{{00L7}{source}{}{sourcepointer-2}{}}
\newlabel{algebra-lemma-filter-minimal-primes-in-support@cref}{{[originallemma][0][]00L7}{[1][0][]source}{}{}{}}
\makeatother
\crefname{originallemma}{original lemma}{original lemmas}
\Crefname{originallemma}{Original lemma}{Original lemmas}
\makeatletter
\newlabel{lemma-symbol-short-exact-sequence}{{02Q3}{source}{}{sourcepointer-3}{}}
\newlabel{lemma-symbol-short-exact-sequence@cref}{{[originallemma][0][]02Q3}{[1][0][]source}{}{}{}}
\makeatother
\crefname{originallemma}{original lemma}{original lemmas}
\Crefname{originallemma}{Original lemma}{Original lemmas}
\makeatletter
\newlabel{lemma-symbol-compare-modules}{{02Q4}{source}{}{sourcepointer-4}{}}
\newlabel{lemma-symbol-compare-modules@cref}{{[originallemma][0][]02Q4}{[1][0][]source}{}{}{}}
\makeatother
\title{The symbol of a one dimensional module is the product of the symbols of its minimal prime quotients raised to the local multiplicities}
\author{The Stacks Project Authors}
\date{stacks-project:chow.tex:local-snapshot-20261012; GNU FDL 1.2 or later}
\begin{document}
\maketitle
\noindent\textit{Editorial note.} All delivered mathematical text is the author's own, copied from the pinned Stacks Project chapter file chow.tex; no statement or proof step was rewritten, reordered or removed.\par\medskip
\noindent\textit{Editorial note.} Preamble provenance: the pinned chapter file opens with an input of the shared preamble and carries no preamble text of its own, so the project's single Stacks-wide preamble staged at sources/stacks-chow/preamble.tex is excerpted verbatim as two spans, lines 14-18 (Xy-pic setup) and lines 154-217 (the native theorem declarations and mathematical macros). Nothing was added, removed or redefined; the wrapper class and the amsmath/amssymb/amsthm/graphicx base packages are supplied by the assembler.\par\medskip
\noindent\textit{Editorial note.} The counted statement and proof contain no cross reference, no citation and no included file, so the delivered body is self contained.\par\medskip
\noindent\textit{Editorial note.} Canonical location of the counted proposition: Stacks tag 02Q5 (sources/stacks-chow/pinned-original-source.tex line 17577).\par\medskip
\noindent\textit{Allowed standard prerequisites (editorial).}\begin{itemize}
\item Existence of a filtration of a finite module over a Noetherian ring with prime quotients. Algebra, Lemma labelled algebra-lemma-filter-Noetherian-module (Stacks tag 00L0), invoked at proof lines 17598-17600. Supplies the filtration whose successive quotients are A/p\_j, which is the backbone of the induction.
\item The number of times a minimal prime of the support occurs in such a filtration equals the length of the localisation. Algebra, Lemma labelled algebra-lemma-filter-minimal-primes-in-support (Stacks tag 00L7), invoked at proof lines 17604-17606. Supplies the identification of the multiplicity with the count of occurrences, which is what makes the exponent in the statement appear.
\item Multiplicativity of the symbol in a short exact sequence and comparison across a quotient. Same chapter, Lemmas labelled lemma-symbol-short-exact-sequence (source line 17503) and lemma-symbol-compare-modules (source line 17533), invoked at proof lines 17619-17620. Supply the two recurrence steps of the induction, for a minimal prime quotient and for the residue field quotient respectively.
\end{itemize}\medskip
\ifdefined\hypertarget\hypertarget{sourcepointer-1}{}\fi
\noindent\textit{Source reference.} Original lemma 00L0: sources/stacks-chow/tags.txt, source lines 770--770. Cross-chapter standard prerequisite located in the project tag list staged in the local snapshot; invoked at proof line 17600 of the counted proof. It is retained only as a navigation pointer; every load-bearing step is inside the delivered span.\par\medskip
\ifdefined\hypertarget\hypertarget{sourcepointer-2}{}\fi
\noindent\textit{Source reference.} Original lemma 00L7: sources/stacks-chow/tags.txt, source lines 777--777. Cross-chapter standard prerequisite located in the project tag list staged in the local snapshot; invoked at proof line 17606 of the counted proof. It is retained only as a navigation pointer; every load-bearing step is inside the delivered span.\par\medskip
\ifdefined\hypertarget\hypertarget{sourcepointer-3}{}\fi
\noindent\textit{Source reference.} Original lemma 02Q3: sources/stacks-chow/pinned-original-source.tex, source lines 17503--17503. Same-chapter pointer into the author environment lemma at source line 17503; invoked at proof line 17619 of the counted proof. It is retained only as a navigation pointer; every load-bearing step is inside the delivered span.\par\medskip
\ifdefined\hypertarget\hypertarget{sourcepointer-4}{}\fi
\noindent\textit{Source reference.} Original lemma 02Q4: sources/stacks-chow/pinned-original-source.tex, source lines 17533--17533. Same-chapter pointer into the author environment lemma at source line 17533; invoked at proof line 17620 of the counted proof. It is retained only as a navigation pointer; every load-bearing step is inside the delivered span.\par\medskip
\begin{lemma}
\label{lemma-compute-symbol-M}
Let $A$ be a Noetherian local ring.
Let $M$ be a finite $A$-module with $\dim(\text{Supp}(M)) = 1$.
Let $a, b \in A$ nonzerodivisors on $M$.
Let $\mathfrak q_1, \ldots, \mathfrak q_t$ be the minimal
primes in the support of $M$. Then
$$
d_M(a, b)
=
\prod\nolimits_{i = 1, \ldots, t}
d_{A/\mathfrak q_i}(a, b)^{
\text{length}_{A_{\mathfrak q_i}}(M_{\mathfrak q_i})}
$$
as elements of $\kappa^*$.
\end{lemma}

\begin{proof}
Choose a filtration by $A$-submodules
$$
0 = M_0 \subset M_1 \subset \ldots \subset M_n = M
$$
such that each quotient $M_j/M_{j - 1}$ is isomorphic
to $A/\mathfrak p_j$ for some prime ideal $\mathfrak p_j$
of $A$. See Algebra, Lemma \ref{algebra-lemma-filter-Noetherian-module}.
For each $j$ we have either $\mathfrak p_j = \mathfrak q_i$
for some $i$, or $\mathfrak p_j = \mathfrak m_A$. Moreover,
for a fixed $i$, the number of $j$ such that
$\mathfrak p_j = \mathfrak q_i$ is equal to
$\text{length}_{A_{\mathfrak q_i}}(M_{\mathfrak q_i})$ by
Algebra, Lemma \ref{algebra-lemma-filter-minimal-primes-in-support}.
Hence $d_{M_j}(a, b)$ is defined for each $j$ and
$$
d_{M_j}(a, b)
=
\left\{
\begin{matrix}
d_{M_{j - 1}}(a, b) d_{A/\mathfrak q_i}(a, b) &
\text{if} & \mathfrak p_j = \mathfrak q_i \\
d_{M_{j - 1}}(a, b) & \text{if} & \mathfrak p_j = \mathfrak m_A
\end{matrix}
\right.
$$
by Lemma \ref{lemma-symbol-short-exact-sequence} in the first instance
and Lemma \ref{lemma-symbol-compare-modules} in the second. Hence the lemma.
\end{proof}

\end{document}
